\section{Related Work}
\label{sec:related}
% \textbf{VLA and large BC models.} RT-2, OpenVLA, action chunking, diffusion policy, $\pi_0$
\textbf{Vision-language-action models.}
Behavioral cloning from large demonstration datasets has recently emerged as the dominant paradigm for training generalist robot manipulation policies (see e.g.~\citep{rt22023arxiv,kim2024openvla,black2024pi_0,geminirobotics,wu2023unleashing,nvidia2025gr00tn1openfoundation}).
Two critical ingredients that have enabled this success are action chunking~\cite{zhao2023learningfinegrainedbimanualmanipulation}, which predicts multiple actions for sequential open-loop execution, and using expressive output distributions, such as diffusion~\cite{chi2023diffusion} or autoregressive generation~\cite{rt22023arxiv}, that can capture the multimodality inherent in demonstration data.
A further advancement came from using large pretrained vision-language models as a backbone for language-conditioned generalist policies, yielding vision-language-action (VLA) models~\cite{rt22023arxiv,kim2024openvla}.
These models import large web-scale prior knowledge into closed-loop robot policies.
Recent work has combined VLA backbones with chunked action generation, either via diffusion~\cite{black2024pi_0} or autoregressive tokenization~\cite{pertsch2025fast,minivla}, achieving state-of-the-art generalist manipulation.
While these policies exhibit impressive generalization capabilities~\cite{pi05,geminirobotics}, their performance on any given task is ultimately bounded by the quality and coverage of the teleoperation data they were trained on---achieving reliable success on precision-critical tasks remains difficult when demonstrations themselves are noisy or inconsistent.
%%SL.3.17: This paragraph has poor coverage of recent VLA work. Citations are cheap, and citing as many papers in this area as possible builds good will. A lot of the citations are conflicting or from "adjacent" institutions (Stanford, Google, Berkeley). Would be nice to cite much more broadly.
 
\textbf{Real-world reinforcement learning.}
Reinforcement learning offers a natural way to go beyond the performance ceiling of demonstration data: by practicing on a task, the agent can discover faster, more precise, or more robust strategies that were never demonstrated.
In practice, real-world RL for robotics operates under tight sample budgets, since every robot rollout costs time and wear.
Off-policy actor-critic methods (e.g.~\cite{haarnoja2018soft,lillicrap2015continuous,fujimoto2018td3,Abdolmaleki2018MaximumAP}) address this by reusing transitions stored in a replay buffer, and sample efficiency can be pushed further by increasing the update-to-data ratio~\cite{hussing2024dissectingdeeprlhigh}, though regularization may be necessary to avoid instability~\cite{chen2021randomized}.
Crucially, off-policy methods can also incorporate human demonstration data to bootstrap learning (e.g.~\cite{ball2023efficient}), combining the strengths of imitation and RL.
A growing body of work has developed practical recipes for deploying RL on physical robots, including autonomous data-collection pipelines~\cite{zhu2020ingredients}, efficient learning frameworks such as SERL~\cite{luo2024serl,luo2024precise} and RL$^{100}$~\cite{rl100}, and human-in-the-loop variants that allow operators to intervene and provide corrections during autonomous execution~\cite{luo2024precise}.
These systems have demonstrated that off-policy actor-critic methods, combined with demonstrations and human corrections, can solve contact-rich manipulation tasks within hours of robot time.
However, they typically train small policies from scratch atop standard pretrained visual encoders (e.g., ResNet), forgoing the rich behavioral priors available in modern VLA models.
\MethodName{} bridges this gap by using a frozen VLA as both the perceptual backbone and the behavioral prior for a lightweight online RL policy.
 
\textbf{RL fine-tuning of VLA models.}
A rapidly growing line of work studies how to improve a pretrained VLA via RL.
These approaches vary primarily in \emph{what} is updated and \emph{how} the RL signal is incorporated. 
At one end of the spectrum, several methods update the full VLA model.
RECAP~\cite{pi06vla} trains the entire $\pi^*_{0.6}$ model end-to-end using offline RL via advantage-conditioned policy extraction: a distributional value function estimates per-timestep advantages, and the VLA is trained on all collected data---demonstrations, autonomous rollouts, and human interventions---with an optimality indicator that upweights high-advantage actions.
By iterating between on-robot data collection and offline RL updates, RECAP more than doubles throughput on complex long-horizon tasks such as espresso making, laundry folding, and box assembly.
Other works apply proximal policy optimization (PPO) or variants thereof to VLA fine-tuning (e.g.~\cite{Ren2025DPPO,Chen2025piRL,Li2025SimpleVLA_RL}), though on-policy methods are difficult to extend to real-world RL in a sample-efficient and scalable fashion. 
At the other end, lightweight methods avoid updating the full VLA and instead train a small auxiliary module on top of the frozen model.
ConRFT~\cite{chen2025conrft} freezes the VLA encoder and fine-tunes an action head using a consistency-based training objective together with a learned binary reward classifier, but operates on single-step actions without chunking on short-horizon tasks.
Policy Decorator~\cite{pi_dec} learns a residual policy whose output is scaled by a hand-tuned hyperparameter and added to the frozen VLA's prediction, but has been demonstrated only in simulation with high sample requirements (on the order of millions of steps).
Probe-Learn-Distill (PLD)~\cite{xiao2025selfimprovingvisionlanguageactionmodelsdata} pre-trains a critic with Cal-QL~\cite{nakamoto2023cal} on base-policy rollouts and then learns a single-step residual policy on top of the frozen VLA, optionally distilling the result back into the VLA via supervised fine-tuning.
GR-RL~\cite{li2025grrlgoingdexterousprecise} takes a multi-stage approach for specializing a generalist VLA on a long-horizon shoe-lacing task: it first performs offline filtered BC, and then performs online RL by learning a noise predictor that steers the frozen VLA's diffusion process in latent space~\citep{Wagenmaker2025DSRL}.
DSRL~\cite{Wagenmaker2025DSRL} similarly operates in the diffusion noise space, learning a latent policy that modulates the denoising process to steer actions toward high-return regions.
 

\MethodName{} shares with these methods the goal of improving a pretrained VLA without the cost of full-model RL, but differs in several key design choices.
First, \MethodName{} introduces an \emph{RL token}---a compact readout representation trained to compress the VLA's internal embeddings---that serves as the state observation for a lightweight actor-critic, preserving the VLA's pretrained perceptual structure while enabling efficient online learning.
Second, \MethodName{} operates over \emph{chunked actions} aligned with the VLA's native action interface, shortening the effective decision horizon for temporal-difference learning under sparse rewards at high control frequencies---in contrast to single-step methods~\cite{chen2025conrft,pi_dec,xiao2025selfimprovingvisionlanguageactionmodelsdata} that face a much longer credit-assignment problem.
Third, rather than predicting residuals or latent noise, the \MethodName{} actor is directly \emph{conditioned on the VLA's sampled reference action chunk and regularized toward it}, turning online RL into local refinement of a good VLA prior behavior policy rather than unconstrained search or implicit modulation of a diffusion process.
Together, these choices enable sample-efficient online RL on real robots---improving both success rate and execution speed within a few hours of practice.
